\section{Open questions (mandatory 5, from paper re-read)}

The five questions below are genuinely unresolved after re-reading Debreceni et al. (2024) alongside the on-disk replication. Each has concrete, executable next steps; none is a rhetorical hedge.

\begin{enumerate}[leftmargin=*,label=\textbf{Q\arabic*.}]

\item \textbf{PIDE version drift and target provenance.} Which PIDE version (2.x with only $\alpha,\beta$ vs 3.2+ with raw dose/SF pairs) was used, and does the 318-point NB1RGB count come from raw dose/SF or from re-simulating $S = \exp(-(\alpha D + \beta D^2))$ off the tabulated $(\alpha,\beta)$? Paper \S 2.1 is ambiguous. If the SF targets are re-simulated from smooth $(\alpha,\beta)$ with no residual scatter, the RF can trivially interpolate and RMSE $\to 0$; if they are raw experimental SF with real Poisson-plating scatter, the 0.0196 RMSE would still be leakage-suspect but genuinely striking. \textbf{Next steps:} (i) register with GSI PIDE and download 2.2 and 3.2; confirm which yields $N=318$ for NB1RGB. (ii) Report per-experiment $\alpha,\beta$ and per-point $(D, SF, \mathrm{LET}, \mathrm{experiment\_id})$. (iii) Retrain the RF on both variants and report RMSE separately --- the delta quantifies target-variance vs interpolation directly.

\item \textbf{Leave-experiment-out CV (the real generalization test).} How does the RF perform under Leave-Experiment-Out (LEO) or Leave-Cell-Line-Out (LOCLO) CV, rather than random-point Monte-Carlo? Random-point CV lets the RF see other points from the same experiment at training time --- essentially within-curve interpolation --- and does not measure what the paper implicitly claims (prediction of a \emph{new} experiment). \textbf{Next steps:} (i) reimplement with \texttt{GroupKFold} keyed on \texttt{experiment\_id} (51 native groups). (ii) Run \texttt{LeaveOneGroupOut} over cell lines using the multi-cell-line PIDE subset the paper mentions in \S 3.4. (iii) Report $R^2$, RMSE, and per-fold learning curves --- the delta from random-point MC-CV \emph{is} the leakage magnitude, and Q1's finding depends on this.

\item \textbf{Mechanistic baselines: LQM($\alpha(\mathrm{LET}), \beta(\mathrm{LET})$), MKM, RMF.} How does the RF(dose, LET) compare against a fair physical baseline --- LQM with LET-dependent parameters --- rather than the paper's dose-only LQ straw man? Everyone in heavy-ion radiobiology knows $\alpha,\beta$ vary with LET; the Furusawa (2000) NB1RGB curves are decades-old and analytically parameterisable. The RF's ``gain'' is largely re-learning those curves empirically. \textbf{Next steps:} (i) fit $\alpha(\mathrm{LET})$, $\beta(\mathrm{LET})$ as smooth splines (or use Furusawa/Weyrather parameterisations), evaluate LQ($\alpha(\mathrm{LET}), \beta(\mathrm{LET})$) on held-out points. (ii) Add MKM (Kase 2006) and RMF (Carlson 2008) as mechanistic baselines. (iii) Report $R^2$, RMSE, and residual structure side-by-side --- genuine ML advantage would show as systematic residual patterns only the RF captures.

\item \textbf{Uncertainty quantification for RF predictions.} How does the RF quantify prediction uncertainty, and is that uncertainty calibrated against experimental replicate variance? The paper reports point predictions and a single RMSE; for the paper's stated clinical/treatment-planning motivation you need calibrated per-prediction intervals (quantile-regression forests, RF infinitesimal jackknife, or conformal prediction). Uncalibrated predictions at novel ion species or clinical energies could be catastrophically overconfident precisely where training data is sparsest. \textbf{Next steps:} (i) retrain as a Quantile Regression Forest (or use \texttt{forestci} for RF-IJ variance) → 80/95\% prediction intervals per point. (ii) Evaluate coverage stratified by LET bin and ion species. (iii) Compare interval width to experimental replicate scatter in PIDE (where multiple experiments exist for the same (ion, LET, dose)); intervals narrower than replicate scatter are miscalibrated.

\item \textbf{Transfer to novel ion species and clinical energies.} Does the RF transfer to (a) helium, proton end-of-range, oxygen --- ion species outside the C/Ne/Si/Fe training mix --- and (b) clinical energies (100--400 MeV/u for carbon therapy) not densely represented in training? Training mix is heavily $^{12}$C-weighted (24 of 51 experiments) with very few $^{28}$Si (7) and $^{56}$Fe (5). Clinical hadrontherapy uses primarily protons and $^{12}$C, and space-radiation exposure is a GCR mix; extrapolation there is where the model would actually be used and there is no evidence in the paper that it works. \textbf{Next steps:} (i) leave-one-ion-out CV over \{C, Ne, Si, Fe\}; report $R^2$/RMSE per held-out ion. (ii) Evaluate the trained RF on the PIDE proton subset (or PARTNER/NASA GCR data) without retraining --- $R^2$ drop = out-of-distribution failure magnitude. (iii) Compare against measured NB1RGB / CHO-K1 survival at HIMAC / GSI / HIT at clinical 200--400 MeV/u, outside the training LET density peak, and report both extrapolation $R^2$ and calibration.

\end{enumerate}
